\documentclass[12pt]{article}

% House style preamble
\input{.house-style/preamble.tex}

% Update PDF metadata
\hypersetup{
  pdftitle={Form-meaning transparency: compositionality, compounds, and priming},
  pdfkeywords={form-meaning transparency, compositionality, construction grammar, compound transparency, masked priming, projectibility}
}

\title{Form--meaning transparency: compositionality, compounds, and priming}
\author{Brett Reynolds \orcidlink{0000-0003-0073-7195}%
\thanks{Contact: \href{mailto:brett.reynolds@humber.ca}{brett.reynolds@humber.ca}}\\
Humber Polytechnic \& University of Toronto}
\date{\today}

\begin{document}
\maketitle

\begin{abstract}
TODO. This paper develops the form--meaning arm of the mediated-accessibility programme. Three within-domain windows on form--meaning transparency (constructional compositionality, rated compound transparency, and masked-priming morphological transparency) each license within-domain projection. The open question is whether they are joint reflexes of a single underlying organization or three separate profiles; a within-items cross-modal probe is designed to settle it. (Spun out of the syntactic umbrella paper, \enquote{Linguistic transparency as mediated accessibility}; this draft is a seed, see \texttt{paper/notes/migration-notes.md}.)
\end{abstract}

\noindent\textbf{Keywords:} form-meaning transparency; compositionality; construction grammar; compound transparency; masked priming; projectibility

\section{Introduction}\label{sec:fm-intro}

TODO: write the standalone introduction. The companion paper (\enquote{Linguistic transparency as mediated accessibility}) treats transparency as a family of mediated-accessibility relations under the schema \mention{$X$ is transparent with respect to $P$ for $R$ when $P$, which $X$ might otherwise have concealed from $R$, remains accessible through $X$ for the purposes of $R$}, and delivers the syntactic profiles (number-transparent quantificational nouns, transparent free relatives). This paper takes up the form--meaning arm in its own right: the relation $R$ is interpretation, learning, or processing rather than agreement, and the question is whether the three windows below cohere into a cross-domain projectibility cluster.

% --- Seed body: sections migrated verbatim from the umbrella paper (2026-06-16). ---
% Cross-references to the umbrella's sections (sec:1, sec:2*, sec:6, sec:7) are currently
% dangling and print as "??"; the QN-revisited (sec:3-4) and TFR-revisited (sec:3-5)
% subsections and the referential-boundary subsection (sec:4-3) belong to the umbrella and
% should be trimmed here. See paper/notes/migration-notes.md before building.
\input{paper/sections/03-constructional-transparency.tex}
\input{paper/sections/04-semantic-transparency.tex}
\input{paper/sections/05-processing-transparency.tex}

\appendix
\input{paper/sections/B1-appendix-b.tex}
\input{paper/sections/C1-appendix-c.tex}
\input{paper/sections/D1-appendix-d.tex}

\section*{Acknowledgements}

% TODO: Replace [MODEL LIST] with the actual models used on this project and their current versions.
The large language models [MODEL LIST] served as drafting and editing aids throughout the preparation of this paper. I am responsible for all theoretical claims, arguments, errors, and interpretive choices.

\clearpage
\printbibliography

\end{document}
